\chapter{Conclusiones y recomendaciones}
\section{Resultados del proyecto académico}
\begin{enumerate}
\item Se desarrolla un proyecto académico de diseño elaborado por estudiantes, con operación aislada. No se declara expediente constructivo ni permisos obtenidos.
\item Los antecedentes de 200 kW, 50 m y Banki siguen sin placa o medición independiente. Las casas de máquinas candidatas difieren 152 m. El salto bruto de gabinete \dato{Hb} m requiere nivelación.
\item PISCO 1981--2025 permite una transposición exploratoria por área geodésica, no un aforo oficial. La estabilidad geométrica 30/90 m no valida el recurso; el control de salida D8 falló. Se requiere calibración local.
\item Limbani y tres registros de Huancasayani suman \dato{PopBase} habitantes y \dato{VivsBase} viviendas ocupadas en 2017. Se proyectan \dato{Users} conexiones potenciales, sin sumar el distrito. Las bases rurales de 2023 producen demanda de \dato{Demand} kW.
\item Se compararon 150, 200, 400 y 600 kW. Para estudiar cobertura del escenario conservador se dimensionaron tres grupos hidráulicos de 200 kW, con caudal nominal \dato{Qnom} \(\mathrm{m^3/s}\), salto neto \dato{Hn} m y tubería nominal de 1000 mm por 350 m. Entregan aproximadamente \dato{Pusers} kW a usuarios en el punto nominal; su margen es reducido y requiere datos garantizados.
\item La potencia hidroeléctrica disponible no es permanente: el indicador diario de excedencia de 95\% del modelo es \dato{Pfirmstudy} kW en usuarios y el mínimo es nulo. Por ello no se garantiza suministro continuo de 150--200 kW en estiaje. La energía hidro servida es \dato{EHidro} MWh/año y el déficit \dato{Edef} MWh/año bajo el perfil supuesto. Se estudia respaldo modular y gestión de cargas.
\item El MOC exige coordinar ELC y cierre lento: a ocho segundos se estiman \dato{PNslow} bar. El primer pico brusco alcanza \dato{PNfast} bar y puede cavitar. PN40 es candidato, sin certificación contra vacío o colapso de cavidades.
\item El sistema eléctrico inicial es aislado, 380 V / 10 kV, con transformador principal de 800 kVA, seis subestaciones y auxiliares con neutro propio. La caída MT de estudio es \dato{DropMT}\%; la rama BT ejemplar tiene \dato{DropBT}\%. La corriente inicial máxima de barra estimada es \dato{IKmax} kA. La protección debe reconocer el decremento de falla de generadores; sobrecorriente simple y una resistencia de tierra aislada no bastan para demostrar seguridad.
\item La matriz separa obligaciones peruanas y recomendaciones. Se actualizan bases rurales 2023 y E.030 de 2026. Los 600 kW superan el umbral de concesión hidráulica de 500 kW; agua, ambiente y obras requieren procedimientos propios.
\item El presupuesto estimado de S/\dato{CapexM} millones presenta escenarios financieros desfavorables. Los 600 kW cubren mejor el pico, pero no demuestran inversión óptima. Se deben comparar rehabilitación y etapas con usuarios y recurso verificados.
\end{enumerate}

\section{Prioridades para continuar}
Antes de escoger capacidad se requieren aforos en captación, nivelación y padrón con demanda medida. Después corresponden crecidas, geotecnia, resistividad, curvas de fabricante, cotizaciones y permisos. Estos datos permitirán comparar rehabilitación, ampliación por etapas y otras fuentes. La eventual conexión al SEIN exige estudios y aprobación propios.
